\subsection{Global coupling of the orthogonal blocks}

\begin{statement}{D23.P20}{Proved}{Innerness of the normalized Hurwitz action}
\end{statement}
Let $\Gamma$ be the degree-1344 permutation image of
D23.P17, and let $N$ be the kernel of its degree-12 block
action. Use the group structure proved in D23.F30:
\[
 H=PSO^-_6(3),\quad S=P\Omega^-_6(3),\quad
 N'=S^{12},\quad N/S^{12}=\mathcal C\subset\mathbb F_2^{12}.
\]
Here $\mathcal C$ consists exactly of the binary vectors whose
sums on the four parallel triples
\[
 \{0,9,11\},\quad\{1,6,7\},\quad
 \{2,4,10\},\quad\{3,5,8\}                             \tag{D23.P20.1}
\]
are equal. Thus $\dim\mathcal C=9$.

For an actual multiplier-one full-GT element with
$\rho\equiv0\pmod3$, the permutation $V_h$ of D23.P18
belongs to $N$. In particular, its automorphism of
$\Gamma$ is inner. This conclusion applies to the normalized
mixed state $(\rho,a_3,a_2)=(0,1,0)$ obtained in D23.H2.

On the four blocks $7,8,10,11$, the three local possibilities
per block in D23.H3 must have indices
\[
                  (i,i,i,j),\qquad 0\leq i,j<3.            \tag{D23.P20.2}
\]
Thus nine assignments remain on those four blocks, for each
of the two cubic bits. This is only a necessary finite
constraint; no actual mixed full-GT element is constructed.

\begin{proof}
We first recall why the actual permutation fixes the twelve
blocks. The first-pair block quotient is the marked
$SL_2(\mathbb F_3)$ quotient and has $f$-value $M_\rho$.
The mixed block commutator quotient is the $C_3$ quotient
of the marked $A_4=S_4'$ value; the supplied FP.37 formula
$f(p,t)=t^{2a_2}p^\rho$ has trivial image in it when
$\rho=0$. These exact marked quotient identifications are
recorded in
\path{hurwitz_cubic216_block_character_projections.json}.
Thus for $\rho=0$ all three transformed braid generators
have their original degree-12 images. The actual induced
permutation centralizes that degree-12 action. Its
centralizer is trivial: a commuting permutation is determined
by the image of one point, and propagation through the
three generators rejects eleven of the twelve possible
images. The certificate replays this complete calculation.
Hence $V_h$ fixes every block individually.

Choose the twelve braid transversals used in D23.F30, so
that every transition map of an element of $\Gamma$ is an
element of $H$ in these local coordinates. Write $V_j$ for
the local permutation of $V_h$ on block $j$. D23.P19
gives its linear orthogonal realization. The exact exhaustive
lift calculation D23.H3 shows that on each of blocks
$7,8,10,11$, every possibility has determinant one and
multiplier one. Thus at least one $V_j$ belongs to $H$.

For any other block $k$, choose $g\in\Gamma$ carrying
block $j$ to $k$. Because $V_h$ normalizes $\Gamma$ by
D23.P18 and fixes all blocks, the element
$g'=V_hgV_h^{-1}$ has the same source and target blocks.
Let $L,L'\in H$ be their normalized transition maps. Then
\[
                         V_k L=L'V_j,
\]
so $V_k=L'V_jL^{-1}\in H$. This propagates membership in
$H$ to all twelve blocks. No classification of the entire
permutation normalizer of $H$ is needed for this step.

The quotient $H/S$ has order two. Let
$v=(v_0,\ldots,v_{11})\in\mathbb F_2^{12}$ be the vector
of its values on the $V_j$. Since $N$ contains $S^{12}$
and has binary image $\mathcal C$, it is exactly the
preimage of $\mathcal C$ in $H^{12}$. Normalization of
$\Gamma$ implies
\[
                       v+qv\in\mathcal C
                 \quad\hbox{for all }q\in\Gamma/N.       \tag{D23.P20.3}
\]
Indeed $V_hgV_h^{-1}g^{-1}\in N$ has precisely this binary
vector, where $q$ is the block permutation of $g$.

We now prove that (D23.P20.3) forces $v\in\mathcal C$.
Let $\pi:\mathbb F_2^{12}\to\mathbb F_2^4$ take the four
triple sums in (D23.P20.1). Then
\[
 \mathcal C=\pi^{-1}\langle(1,1,1,1)\rangle.
\]
On the four triples the three braid generators induce the
permutations
\[
      (1\,2\,3),\qquad(0\,1\,3),\qquad(1\,2\,3).
\]
They have order three and together act transitively. If the
class of $w=\pi(v)$ modulo the constant vector is fixed by
one such permutation $a$, write
$aw+w=\epsilon(1,1,1,1)$. Adding its three translates gives
$0=3\epsilon(1,1,1,1)=\epsilon(1,1,1,1)$ in
characteristic two. Hence $aw=w$. Applying this to the
three generators makes $w$ a constant vector, by
transitivity. Thus $v\in\mathcal C$, as claimed.
Since $V_h\in H^{12}$, this proves $V_h\in N$.

Finally D23.P18 gives $V_hA=AV_h$ for the first braid
generator $A$. Its block permutation cycles
$7\mapsto10\mapsto8\mapsto7$ and fixes $11$.
The exact local transitions show that this equation
identifies the three H3 candidate indices on the cycled
blocks. All three possibilities on block 11 commute with
its local $A$ transition. Checking the $3^4$ possibilities
therefore gives exactly (D23.P20.2), independently for
$a_3=0$ and $a_3=1$.
\end{proof}

\paragraph{Exact certificate and scope.}
\path{hurwitz_global_innerness.py} replays the degree-12
centralizer, the four-triple action, all 4096 binary sign
vectors, and both local $3^4$ coupling calculations. Its
record is \path{hurwitz_global_innerness.json}. The binary
enumeration agrees with the short proof: exactly the 512
code vectors normalize, and exactly 32 remain after the
four H3 signs are set to zero. This last count uses the
separately checked H3 spinor values; innerness itself uses
only their determinant-one, multiplier-one membership in
$H$.

The group structure input is
\path{five_level_six/hurwitz_block_structure.tex}, D23.F30;
the local exhaustive lift input is D23.H3 with
\path{work/code/hurwitz_h1_orthogonal_lifts.py} and its
JSON/NPZ records. The necessary actual-source interpretation
comes from D23.P18--P19. Innerness of this finite quotient
does not establish innerness of the full profinite GT
action, nor the vanishing of $\Delta$.
